
\subsubsection{2.1 Pretraining Corpus Filtering}

Early corpus filtering relied on perplexity thresholds, selecting documents with low GPT-2 perplexity as proxies for quality \citep{Wenzek2020}. RefinedWeb demonstrated that perplexity filtering combined with deduplication could match human-curated datasets \citep{Penedo2023}.

DataComp-LM introduced classifier-based filtering: train a FastText model to distinguish high-quality from low-quality documents, then filter corpora at scale \citep{Li2024}. Their approach enabled a 7B model trained on 40\% less data to achieve 64\% MMLU accuracy. FineWeb-Edu extended this by training educational quality classifiers, producing +5.0pp MMLU gains \citep{Penedo2024}.

These methods optimize for pretraining objectives measured through language modeling loss and downstream task performance after training. Whether perplexity or educational quality correlate with retrieval utility---where models must locate information at inference time---remains untested. Our work applies model-based filtering methodology to retrieval-specific training signals.

\subsubsection{2.2 Retrieval Benchmarks}

BEIR provides heterogeneous evaluation across 18 datasets spanning factoid QA, argument retrieval, and domain-specific search \citep{Thakur2021}. BEIR's relevance annotations (qrels) indicate which corpus documents answer each query, enabling controlled evaluation.

Dense Passage Retrieval (DPR) uses bi-encoder architectures for dense retrieval, achieving 79.4\% top-20 accuracy on Natural Questions versus BM25's 59.1\% \citep{Karpukhin2020}. However, DPR's success depends on corpus coverage---retrieval cannot succeed if relevant documents are filtered out during corpus construction.

Prior work measures retrieval model quality on fixed corpora. We invert this: use BEIR annotations as training data for corpus filtering, evaluating corpus quality using a fixed retrieval model (DPR).

\subsubsection{2.3 Quality Metrics for Text}

Entity-based metrics have been used heuristically in corpus filtering under the assumption that factual density correlates with informativeness. Wikipedia exhibits high named entity density compared to conversational text (Färber et al., 2018). However, whether entity density correlates with retrieval performance has not been rigorously tested.

Educational quality classifiers use perplexity, document structure, and n-gram statistics \citep{Penedo2024}, but these features optimize for learnability (how well models internalize text during training) rather than retrievability (how well text supports fact lookup at inference time).

We explicitly test the entity density hypothesis by measuring NER-based factual density in classifier-selected corpora and comparing against perplexity-matched baselines.
